\section{A weighted projected impedance estimate}
\label{sec:impedance}

We retain the boundary geometry and the conditional boundary law of
\cref{cell:geometry,cell:branch}. In this section the interior of the cell
does not enter the calculation. Its only relevant feature is that it
couples to the boundary space through the projection $P$. We write
\[
 P=P_0+Q,
 \qquad Q=Q_1+Q_2+Q_3+Q_4,
\]
where $P_0$ projects onto the single side sites and $Q_a$ projects onto the
two-dimensional plane of the $a$th corner triple. The four complementary
directions in the triples are denoted by $n_a$. In the prescribed ordering
of each triple, $n_a=(-1,-1,1)/\sqrt3$. Thus $\rank Q=8$, whereas the
dimension of the whole boundary space is $O(L)$.

All norms below use the fixed real orthonormal bases of these spaces.
Constants may depend on the fixed disorder, profile, and corner-box
parameters. They will not depend on the exterior matrices, their norms,
the source vector, or the artificial-probe mask. Until an exponent $H$
is selected below, every expression denoted by $\operatorname{poly}(L)$
means a bound $C(1+L)^C$ whose constants are independent of $H$. A final
polynomial bound is allowed to depend on that selected, fixed exponent.

\begin{proposition}[Weighted projected impedance]
\label{imp:weighted}
Let $T$ be any fixed complex symmetric matrix on the boundary space with
$-\operatorname{Im}T\ge0$, and let $k$ be any fixed vector there. For each
boundary site prescribe $a_i$ to be one of $0$, $1$, and $\vartheta$, with
the prescription independent of the boundary variables. Set
\begin{equation}
 \mathcal T=T+\diag_B\bigl(v_i-i\tau a_i(v_i)\bigr),
 \qquad
 \rho=\norm{P(\mathcal T-iP)^{-1}k}^2.
 \label{imp:rho}
\end{equation}
Fix $c_*>0$ and suppose $0\le\tau\le e^{-c_*L}$. There are constants
$c>0$, $C<\infty$, and a sufficiently large lower threshold for $L$ such
that
\begin{equation}
 \E_B\!\left[
   \rho\,\one_{\{\norm{P\mathcal T^{-1}P}\le C(1+L)^C\}}
 \right]
 \ge c\,\E_B\rho.
 \label{imp:weighted-estimate}
\end{equation}
Here the expectation uses the boundary law conditional on its four
corner boxes. The constants are uniform in all the exterior data just
described. Moreover, changing only the single side variables to any other
bounded allowed settings, including their artificial absorptions,
changes $\rho$ by at most fixed multiplicative factors, pointwise wherever
the inverses exist.
\end{proposition}

The weight in \cref{imp:rho} matters: a small unweighted exceptional set
could contain nearly all the response of a weakly coupled exterior mode.
We first control the inverse on the corner planes $Q$ while leaving unit
absorption on the single side coordinates $P_0$. The corner estimates
must retain a fixed fraction of the response weight. We then average the
common real side shift to bound the whole projected inverse
$P\mathcal T^{-1}P$ with the auxiliary $P_0$ absorption removed.

\subsection{Elementary reductions and polynomial estimates}

We first record the comparison that permits bounded changes of the
absorbed coordinates.

\begin{lemma}[Comparison of absorbed responses]
\label{imp:comparison}
Suppose $D_1,D_2$ are invertible matrices with
$-\operatorname{Im}D_j\ge P$. If
$D_2-D_1=P\Delta P$ and $\norm{\Delta}\le K$, then, for every $k$,
\[
 (1+K)^{-1}\norm{PD_1^{-1}k}
 \le \norm{PD_2^{-1}k}
 \le (1+K)\norm{PD_1^{-1}k}.
\]
\end{lemma}

\begin{proof}
The inverse difference identity gives
\[
 PD_2^{-1}k
 =PD_1^{-1}k-PD_2^{-1}P\Delta PD_1^{-1}k.
\]
By \cref{mat:ward}, $\norm{PD_2^{-1}P}\le1$. This proves the upper
bound; interchanging the indices proves the lower bound.
\end{proof}

The inverses used in this section exist almost everywhere in the stated
boundary law. To see this without any nonsingularity assumption on $T$,
recall that a kernel vector of a passive matrix is annihilated by both its
Hermitian real and imaginary parts. Absorption on $P$ restricts such a
vector to the four missing directions. On that space the Hermitian real
part has the form of a fixed real symmetric matrix plus the independent
common corner shifts on its diagonal. Its determinant is a nonzero
polynomial, with the product of those shifts as its leading monomial.
The same reasoning applies after one missing direction is deleted or
after bounded changes in plane blocks. With absorption only on $P_0$,
a possible kernel lies on the full corner triples. The four common shifts
act there as scalar shifts on their respective three-dimensional blocks;
the determinant again has a nonzero leading monomial. Finally,
$\mathcal T$ itself has physical real diagonal variables with a joint
density, so its Hermitian real part is nonsingular
almost everywhere. These observations also justify almost-everywhere
slicing below.

The response weights are integrable. For example, writing
$G=(\mathcal T-iP)^{-1}$, the Ward identity implies
\[
 \norm{PGk}^2\le k^*(\operatorname{Im}G)k
 \le \norm{k}^2\tr\operatorname{Im}G.
\]
The joint law of the physical boundary diagonals has a bounded density on
a compact set; the bound here may depend on $L$. Slice that set along a
common real shift of all the physical diagonals. The phase estimate
\cref{mat:phase}, including its artificial-diagonal derivative term,
bounds the integral of $\tr\operatorname{Im}G$ along each slice. Integration
over the transverse coordinates proves finiteness. One may add
$-i\zeta\one$ first and pass to the limit by Fatou. In particular, after
some boundary variables are frozen, all integrals below are finite for
almost every such choice.

We shall use two elementary facts about polynomials. Their degree bounds
are fixed throughout this section.

\begin{lemma}[Polynomial norms and products]
\label{imp:polynomials}
On a rectangular box in a fixed number of real variables, let $p$ be a
polynomial of degree at most $d$ with values in a finite-dimensional
Hilbert space. There is a constant depending only on $d$ and the number
of variables such that
\begin{equation}
 \sup\norm{p}^2\le C\,\mathop{\mathrm{avg}}\norm{p}^2.
 \label{imp:poly-average}
\end{equation}
If the length of a coordinate interval is at least $\ell$, then
\begin{equation}
 \sup\norm{\partial_jp}\le C\ell^{-1}\sup\norm p.
 \label{imp:poly-derivative}
\end{equation}
If $q$ is a scalar polynomial of degree at most $d$, then
\begin{equation}
 \sup\norm{pq}\ge c\,\sup\norm p\,\sup\abs q.
 \label{imp:poly-product}
\end{equation}
The constants do not depend on the dimension of the target Hilbert space
or the side lengths of the box.
\end{lemma}

\begin{proof}
Affinely rescale the box to a unit cube. For scalar polynomials of bounded
degree, the asserted norm comparisons and derivative estimate follow
from equivalence of norms on a fixed finite-dimensional vector space.
For \cref{imp:poly-average}, choose a point where $\norm p$ is maximal
and a unit linear functional attaining that norm at that point; apply
the scalar estimate to its composition with $p$. The derivative estimate
follows in the same way by projecting at a maximizing derivative.

For the scalar product estimate, normalize the two supremum norms to
one. The two unit spheres of bounded-degree scalar polynomials are
compact. The supremum norm of their product has a positive minimum:
otherwise two nonzero limiting polynomials would have identically zero
product. Projecting the vector polynomial at one of its maximizing
points reduces \cref{imp:poly-product} to this scalar result.
\end{proof}

The following one-dimensional observation will prevent an extremely
narrow response peak from defeating a short-exception estimate.

\begin{lemma}[Resolution of a weighted slice]
\label{imp:resolution}
Let $I$ be an interval of length at least $\ell$, let $N$ be a
Hilbert-space-valued polynomial of fixed degree on $I$, and let $z$ be a
continuously differentiable complex function satisfying
\[
 \abs{z(s)}\ge c_0 a,
 \qquad \abs{z'(s)}\le C_0,
 \qquad a>0.
\]
Then
\[
 \sup_{s\in I}\frac{\norm{N(s)}^2}{\abs{z(s)}^2}
 \le \frac{C}{\min(\ell,a)}
       \int_I\frac{\norm{N(s)}^2}{\abs{z(s)}^2}\,\dd s.
\]
Here $C$ depends only on $c_0,C_0$ and the degree bound.
\end{lemma}

\begin{proof}
Take a maximizing point $s_0$ and an interval $I_0\subset I$ containing it,
of length a sufficiently small fixed multiple of $\min(\ell,a)$.
The derivative bound gives $\abs{z(s)}\le2\abs{z(s_0)}$ on $I_0$.
Applying \cref{imp:poly-average} on $I_0$ yields
\[
 \int_{I_0}\frac{\norm{N(s)}^2}{\abs{z(s)}^2}\,\dd s
 \ge \frac{c\abs{I_0}\norm{N(s_0)}^2}{\abs{z(s_0)}^2}.
\]
This proves the claim, including when $s_0$ is an endpoint.
\end{proof}

\subsection{Eight corner tests}

Remove the common shift from the single side variables and omit their
artificial absorptions. Denote the resulting boundary matrix by
$\mathcal T_0$, and set
\[
 w_0=\norm{P(\mathcal T_0-iP)^{-1}k}^2.
\]
The independent single side uniforms are still present. By
\cref{imp:comparison}, $w_0$ and $\rho$ are pointwise comparable by
fixed factors. We now hold those independent side uniforms fixed.

Choose an exponent $H$, to be fixed after the polynomial losses below,
and put
\[
 \eta=L^{-H},
 \qquad G_\eta=(\mathcal T_0-iP_0-i\eta Q)^{-1}.
\]
For each vector $b'$ in the eight prescribed real orthonormal corner-plane
bases, define
\begin{equation}
 \mathcal E_{b'}=
 \left\{\eta\operatorname{Im}(b'^\top G_\eta b')\ge\frac1{40}\right\}.
 \label{imp:test}
\end{equation}
Outside their union,
\begin{equation}
 \norm{Q(\mathcal T_0-iP_0)^{-1}Q}\le\eta^{-1}.
 \label{imp:Qgood}
\end{equation}
Indeed, if $\norm{QG_\eta Q}<(2\eta)^{-1}$, inverse difference and a
Neumann series give the bound in \cref{imp:Qgood}. Hence failure of that
bound forces $\norm{QG_\eta Q}\ge(2\eta)^{-1}$. Ward on $Q$ then gives
\[
 \eta\tr(Q\operatorname{Im}G_\eta Q)
 \ge \eta^2\norm{QG_\eta Q}_{\mathrm{HS}}^2\ge\frac14.
\]
If all eight tests fail to occur, the left side is less than $8/40<1/4$,
a contradiction.

Fix one test and condition on all boundary variables outside its corner
triple. Use centered box coordinates $(s,t_1,t_2)$ there, where $s$ is the
common shift and $t_j$ is the shift whose diagonal matrix maps the missing
direction $n$ to the plane basis vector $r_j$. In particular, the $t_j$
shift has zero $n,n$ entry. Write $t=(t_1,t_2)$. Every side length of the
box is at least $\ell\ge(1+L)^{-C}$ and at most a fixed constant. All
probabilities and averages on this box are normalized uniform ones.

There is first an unweighted bound
\begin{equation}
 \PP_{s,t}(\mathcal E_{b'})\le\operatorname{poly}(L)\eta.
 \label{imp:unweighted}
\end{equation}
For fixed $t$, apply \cref{mat:phase} while varying $s$. Its real
derivative is the identity on the tested triple. The artificial
derivatives satisfy the profile condition and are controlled by their
actual absorption, so the integral of the triple partial trace of
$\operatorname{Im}G_\eta$ is polynomially bounded. Since $b'$ is a unit
vector in that triple, its quadratic form is bounded by this trace.
Divide by the interval length, then average over $t$, to obtain
\cref{imp:unweighted}. Neither the phase range nor this polynomial loss
depends on $H$ or on the exterior norm.

We must strengthen \cref{imp:unweighted} to a bound for $w_0$-mass.

\subsection{A fixed remainder at one corner}

The exceptional test still uses the actual inverse $G_\eta$. To measure
its weight, we replace $w_0$ by a comparable response whose remainder
after deleting the tested missing direction $n$ is fixed. In
$\mathcal T_0-iP$, freeze the real variations of the tested triple
within its plane-plane block at the box center. Replace the artificial
absorption in that plane-plane block by $\tau$ times the plane identity.
Retain every fixed term from $T$ and retain all actual mixed terms and
the missing-direction diagonal. Call this modified matrix
$\widetilde D$.

The modification is bounded and supported on the tested plane. It also
preserves absorption at least $P$ plus the original artificial absorption
of the tested triple. In fact the change of artificial absorption is
\[
 \tau Q_a\bigl(\one-\diag(a_i(v_i))\bigr)Q_a\ge0,
\]
because $0\le a_i\le1$. Consequently
\begin{equation}
 C^{-1}w_0\le W:=\norm{P\widetilde D^{-1}k}^2\le Cw_0.
 \label{imp:modified-weight}
\end{equation}
The constant is fixed by \cref{imp:comparison}.

Delete the tested missing direction $n$, and denote the inverse of the
remaining block of $\widetilde D$ by $G_{\mathrm a}$. This inverse is
independent of $(s,t)$. The off-diagonal column from $n$ and its Schur
denominator have the form
\begin{align}
 p_1(s,t)&=p_{10}+t_1r_1+t_2r_2
       -i\tau\sum_{i\text{ in the triple}}a_i(v_i)n_iQe_i,
 \label{imp:p1}\\
 s-e(s,t)&=s-e_{\mathrm{fix}}
       -p_1^\top G_{\mathrm a}p_1
       -i\tau\sum_{i\text{ in the triple}}a_i(v_i)n_i^2.
 \label{imp:denominator}
\end{align}
The fixed vector $p_{10}$ may have components outside the tested triple;
no bound on it is assumed. All variable parts of $p_1$ lie in the tested
plane. The constants in \cref{imp:denominator} include the fixed diagonal
at $n$; there is no real $t$-dependent contribution to that diagonal.

Write $k_{\widehat n}$ for $k$ with its $n$ coordinate deleted, and set
\begin{equation}
 a=PG_{\mathrm a}k_{\widehat n},\qquad
 d_1=PG_{\mathrm a}p_1,\qquad
 \beta=p_1^\top G_{\mathrm a}k_{\widehat n}-k_n.
 \label{imp:adb}
\end{equation}
The block inverse formula yields
\begin{equation}
 W=\norm{a+\frac{d_1\beta}{s-e}}^2
   =\frac{\norm{N_1}^2}{\abs{s-e}^2},
 \qquad N_1=a(s-e)+d_1\beta.
 \label{imp:W}
\end{equation}
The norm of $d_1$ measures the response on the absorbed coordinates
induced by unit amplitude in the missing direction. The scalar Schur
denominator $s-e$ controls how sharply $W$ can concentrate. It need not agree with
the corresponding denominator in the actual test $G_\eta$.
Each profile lies on one fixed polynomial piece throughout the corner
box. Hence $p_1,d_1,\beta,e,N_1$ are polynomials in $(s,t)$ of bounded
degree, independently of all exterior data. Their coefficients need not
be bounded.

\begin{lemma}[Corner Schur estimates]
\label{imp:schur-estimates}
Uniformly throughout the box,
\begin{align}
 \operatorname{Im}e&\ge\norm{d_1}^2,
 &\abs{\partial_s e}&\le C\tau(1+\norm{d_1}),
 \label{imp:derivatives-a}\\
 \norm{\partial_s d_1}&\le C\tau,
 &\abs{\partial_s\beta}&\le C\tau\norm a,
 \label{imp:derivatives-b}\\
 \norm{\partial_{t_j}d_1}&\le C,
 &\abs{\partial_{t_j}e}&\le C(1+\norm{d_1}).
 \label{imp:derivatives-c}
\end{align}
If at least one artificial probe is present in the tested triple and
$\tau>0$, then also
\begin{equation}
 \operatorname{Im}e\ge\frac{\tau}{\operatorname{poly}(L)}.
 \label{imp:probe-floor}
\end{equation}
For every fixed $t$,
\begin{equation}
 \int\frac{\norm{d_1(s,t)}^2}{\abs{s-e(s,t)}^2}\,\dd s
 \le\int\operatorname{Im}\frac1{s-e(s,t)}\,\dd s
 \le\operatorname{poly}(L).
 \label{imp:s-integral}
\end{equation}
\end{lemma}

\begin{proof}
Apply absorption to the Schur extension
$(-G_{\mathrm a}p_1,1)$ at $n$. Its quadratic form under $\widetilde D$
is the scalar $s-e$, and its absorbed $P$-coordinates have norm
$\norm{d_1}$. This proves the first inequality. Every derivative of
$p_1$ is supported in the tested plane. Its $s$ derivative has size
$O(\tau)$, whereas its $t_j$ derivative is $r_j+O(\tau)$. Since
$\norm{PG_{\mathrm a}P}\le1$, these facts give the bounds on $d_1$.
Differentiating \cref{imp:denominator,imp:adb}, using transpose symmetry,
gives the other derivative bounds: for example,
\[
 \partial_s e
 =2(\partial_s p_1)^\top G_{\mathrm a}p_1
   +i\tau\sum_i\partial_s(a_i(v_i))n_i^2.
\]
The projection of $G_{\mathrm a}p_1$ onto the support of
$\partial_s p_1$ is controlled by $d_1$. Thus arbitrarily large fixed
components of $p_{10}$ cause no loss.

For \cref{imp:probe-floor}, an active profile in the triple is at least
an inverse polynomial, by the box clearance from the support endpoints.
If $\norm{d_1}\ge1/(2\sqrt3)$, the first bound already supplies a fixed
positive lower bound. Otherwise the corresponding physical coordinate
of $(-G_{\mathrm a}p_1,1)$ has modulus at least $1/(2\sqrt3)$, since
$\abs{n_i}=1/\sqrt3$. Its artificial absorption supplies
$\tau/\operatorname{poly}(L)$.

For the integral bound, let $y=\operatorname{Im}e$ and
$V=\operatorname{Im}(s-e)^{-1}=y/\abs{s-e}^2$. The first inequality
of \cref{imp:s-integral} follows from $y\ge\norm{d_1}^2$. If $e$ is
independent of $s$, ordinary scalar shift integration gives
$\int V\le\pi$; real poles are null sets. Otherwise $\tau>0$ and a
probe is present. Combining \cref{imp:derivatives-a,imp:probe-floor},
\[
 \frac{\abs{\partial_s e}}{\sqrt y}
 \le C\tau\left(\frac1{\sqrt y}+1\right)
 \le\operatorname{poly}(L).
\]
The argument of $s-e$ lies in $(-\pi,0)$, and its derivative is
\[
 \frac{\dd}{\dd s}\arg(s-e)
 =V-\operatorname{Im}\frac{\partial_s e}{s-e}.
\]
The error is bounded by $\operatorname{poly}(L)\sqrt V$. Young's
inequality absorbs half the integral of $V$; the remaining constant and
the phase variation are polynomially bounded on the bounded interval.
This proves \cref{imp:s-integral}.
\end{proof}

\subsection{Weighted exceptions when the coupling is appreciable}

Put
\[
 D=\max_{s,t}\norm{d_1(s,t)},\qquad \alpha_0=\frac1{100}.
\]
All maxima are taken over the closed box, which does not affect its
conditional law. The bound $\operatorname{Im}e\ge\norm{d_1}^2$ ties the
width of the response to this coupling. When $D\ge\eta^{\alpha_0}$, we
can isolate the few transverse parameters where the coupling becomes
small and use the unweighted test bound elsewhere. The case where $D$
is small throughout the box will instead require comparing the length
of each exceptional slice with its response width.

We first treat $D\ge\eta^{\alpha_0}$.

Suppose initially that $\abs{s-e}\ge1$ throughout the box. The maximum
and minimum of this denominator are then comparable by a fixed factor.
Indeed, \cref{imp:derivatives-b,imp:derivatives-c} show that $d_1$ varies
by at most a fixed constant. If $D$ is large, then
$\operatorname{Im}e\ge cD^2$ throughout, whereas the variation of
$s-e$ is at most $C(1+D)$. If $D$ is bounded, the same variation bound
and the assumed lower bound one suffice. By
\cref{imp:poly-average,imp:unweighted},
\begin{equation}
 \E_{s,t}[W\one_{\mathcal E_{b'}}]
 \le\operatorname{poly}(L)\eta\,\E_{s,t}W.
 \label{imp:comparable-case}
\end{equation}

It remains to consider the case where $\abs{s-e}<1$ somewhere. At that
point $\norm{d_1}<1$; the derivative bounds then imply that $D$ is
bounded throughout the box. Since $s$ is bounded, the same argument
shows that $e$ and $s-e$ are bounded as well. Set
\[
 M_1=\max\norm{N_1},\qquad B_1=\max\abs\beta.
\]
The bounded denominator and \cref{imp:poly-average} give
\begin{equation}
 \E_{s,t}W\ge cM_1^2.
 \label{imp:weight-lower}
\end{equation}
Furthermore,
\begin{equation}
 \norm a+DB_1\le\operatorname{poly}(L)M_1.
 \label{imp:numerator-control}
\end{equation}
Here is the required quantitative argument. By the product estimate,
\[
 DB_1\le C\max\norm{d_1\beta}
       \le C(M_1+\norm a).
\]
Differentiate $N_1=a(s-e)+d_1\beta$ in $s$ and use
\cref{imp:poly-derivative,imp:derivatives-a,imp:derivatives-b}. As $D$
is bounded, this gives
\[
 \norm a\le C\ell^{-1}M_1+C\tau(\norm a+B_1).
\]
Substitute the preceding estimate for $B_1$. Since
$D\ge\eta^{\alpha_0}$, the coefficient
$C\tau(1+D^{-1})$ is less than one half for sufficiently large $L$
after $H$ is fixed. Absorption proves \cref{imp:numerator-control}.
Its polynomial exponent is independent of $H$.

Choose $\Delta=\eta^{1/8}$ and let
\[
 \mathcal A=\{t:\min_s\norm{d_1(s,t)}<\Delta\}.
\]
Then
\begin{equation}
 \PP_t(\mathcal A)\le C\frac{\Delta}{D}.
 \label{imp:affine-sublevel}
\end{equation}
To prove this, \cref{imp:p1} and Ward show that $d_1(s,t)$ is within
$C\tau$ of an affine function of $t$, uniformly in $s$. Rescale the two
$t$ intervals to unit intervals. If $\mathcal A$ is nonempty, the affine
function has a value of norm at most $\Delta+C\tau$, and its maximum
norm is at least $D-C\tau$. Project it by a real unit linear functional
attaining that latter norm. Because $\Delta/D\to0$ and $\tau\ll\Delta$,
one of the two scalar slopes has magnitude at least $cD$. Slicing in
that coordinate bounds the fraction where the scalar magnitude is at
most $\Delta+C\tau$ by $C\Delta/D$. This includes $\mathcal A$.

For each $t\in\mathcal A$, \cref{imp:W,imp:s-integral} gives
\[
 \E_s W\le\operatorname{poly}(L)(\norm a^2+B_1^2)
 \le\operatorname{poly}(L)(1+D^{-2})M_1^2.
\]
On the complement of $\mathcal A$, absorption gives
$\abs{s-e}\ge\norm{d_1}^2\ge\Delta^2$, so
$W\le\Delta^{-4}M_1^2$. Using \cref{imp:unweighted} there and
\cref{imp:affine-sublevel} on $\mathcal A$, we obtain
\begin{equation}
 \E_{s,t}[W\one_{\mathcal E_{b'}}]
 \le\operatorname{poly}(L)
     \left(\frac{\Delta}{D^3}+\frac{\eta}{\Delta^4}\right)
     \E_{s,t}W.
 \label{imp:first-case}
\end{equation}
We used the boundedness of $D$ and \cref{imp:weight-lower} in the last
step. If $M_1=0$, all the weighted integrals vanish and the inequality
holds directly. In particular, this argument never divides by a possibly
zero expected weight. The powers in \cref{imp:first-case} satisfy
\[
 \frac{\Delta}{D^3}\le\eta^{1/8-3\alpha_0},
 \qquad \frac\eta{\Delta^4}=\eta^{1/2},
 \qquad \frac18-3\alpha_0>0.
\]

\subsection{Weighted exceptions when the coupling is small}

We now assume $D<\eta^{\alpha_0}$ and write $b'=r_j$ for the tested
basis vector. Vary $t_j$ between two points separated by at least half
its interval length. From \cref{imp:p1,imp:adb},
\begin{equation}
 a_{b'}:=PG_{\mathrm a}b',
 \qquad
 \norm{a_{b'}}\le C\ell^{-1}(D+\tau).
 \label{imp:small-b}
\end{equation}
This can be made as small as a prescribed fixed constant by choosing
$H$ sufficiently large, then $L$ sufficiently large. The inference uses
variation over the whole $t_j$ interval; it does not assume that the
exterior has a lower bound on its coupling.

Next delete $n$ in the actual matrix $G_\eta^{-1}$, without freezing its
plane block. Let $G_{\eta,\mathrm{rem}}$ denote the remaining inverse,
put
\[
 d_{1,\eta}=PG_{\eta,\mathrm{rem}}p_1,
\]
and write $s-e_\eta$ for its Schur denominator at $n$. Its off-diagonal
column is exactly the same $p_1$ as before. The remaining matrices for
$G_{\mathrm a}$ and $G_{\eta,\mathrm{rem}}$ differ by a bounded matrix
supported on $Q$. Moreover,
$\norm{QG_{\eta,\mathrm{rem}}Q}\le\eta^{-1}$ by Ward. Inverse
difference and symmetry therefore give
\begin{align}
 \eta\abs{b'^\top G_{\eta,\mathrm{rem}}b'}
   &\le C\norm{a_{b'}},
 \label{imp:small-comparison-a}\\
 \abs{b'^\top d_{1,\eta}}
   &\le\abs{b'^\top d_1}
       +C\norm{a_{b'}}\norm{Qd_{1,\eta}}.
 \label{imp:small-comparison-b}
\end{align}
For clarity, the inverse identity used in both estimates is
$G_{\eta,\mathrm{rem}}
=G_{\mathrm a}+G_{\mathrm a}\Delta G_{\eta,\mathrm{rem}}$,
where $\Delta=Q\Delta Q$ has bounded norm. The first left factor is thus
controlled by $PG_{\mathrm a}b'$, regardless of the exterior norm.

Take the bound in \cref{imp:small-b} small enough that the first term
in the Schur inverse formula cannot by itself cause
$\mathcal E_{b'}$. On that exception,
\begin{equation}
 \frac{\eta\abs{b'^\top d_{1,\eta}}^2}{\abs{s-e_\eta}}
 \ge c.
 \label{imp:exception-schur}
\end{equation}
On the other hand, absorption on the actual Schur extension yields
\[
 \operatorname{Im}e_\eta\ge\eta\norm{Qd_{1,\eta}}^2.
\]
Together with \cref{imp:exception-schur}, this implies
$\norm{Qd_{1,\eta}}\le C\abs{b'^\top d_{1,\eta}}$. Insert
\cref{imp:small-comparison-b} and absorb its small last coefficient.
We conclude that, on the exception,
\begin{equation}
 \norm{Qd_{1,\eta}}\le C\norm{d_1},
 \qquad
 \abs{s-e_\eta}\le C\eta\norm{d_1}^2.
 \label{imp:small-exception}
\end{equation}

The exceptional set is short on every almost-everywhere $s$ slice.
To establish this rigorously, let
$D_t=\max_s\norm{d_1(s,t)}$. Differentiating the actual Schur formula
shows that, on the exception,
\[
 \abs{\partial_s e_\eta}
 \le C\norm{Qd_{1,\eta}}^2
       +C\tau(1+\norm{Qd_{1,\eta}})
 \le C D^2+C\tau(1+D)<\frac12.
\]
The first term comes from the real derivative of the tested plane
block. Its artificial derivative is absorbed in the same bound; the
mixed and missing-direction derivatives give the terms containing
$\tau$. Thus $\operatorname{Re}(s-e_\eta)$ has derivative at least
one half on the exception.

There are only boundedly many components of that exceptional slice away
from poles. Indeed, with all other variables fixed, the matrix varies
polynomially only on the three coordinates of the tested triple.
In a basis containing those coordinates, expansion of the determinant
uses variable factors from at most three rows. Its degree, and the
degrees of the relevant cofactors, are bounded by a constant depending
only on the profile degree, not on the dimension or norm of the exterior.
The same assertion holds after deleting $n$. Clearing the squared
absolute value of a denominator converts \cref{imp:test} into a real
polynomial inequality of bounded degree. Such a set has a bounded
number of intervals and isolated points, except for poles, which also
have bounded number on an almost-everywhere slice.

On each interval component, \cref{imp:small-exception} confines
$\operatorname{Re}(s-e_\eta)$ to an interval of length
$C\eta D_t^2$, while its derivative is at least one half. Summing over
the bounded number of components gives
\begin{equation}
 \abs{\{s:(s,t)\in\mathcal E_{b'}\}}
 \le C\eta D_t^2.
 \label{imp:slice-length}
\end{equation}
If $D_t=0$, \cref{imp:small-exception} and the Schur formula rule out
the exception off its null set of poles.

For $D_t>0$, the denominator of the modified weight has a matching
minimum scale:
\begin{equation}
 \operatorname{Im}e(s,t)\ge cD_t^2
 \quad\hbox{throughout the slice},
 \qquad
 \abs{\partial_s(s-e)}\le C.
 \label{imp:linewidth}
\end{equation}
Here is the point requiring the artificial-absorption floor. By
\cref{imp:derivatives-b}, the oscillation of $d_1$ on the slice is
$O(\tau)$. If its norm is everywhere at least $D_t/2$, absorption proves
the first bound immediately. Otherwise $D_t\le C\tau$. In this latter
case an artificial term must be present, since without one $d_1$ is
independent of $s$. Now \cref{imp:probe-floor} gives
$\operatorname{Im}e\ge\tau/\operatorname{poly}(L)\ge cD_t^2$
for sufficiently large $L$. The derivative bound follows from
\cref{imp:derivatives-a} and $D<\eta^{\alpha_0}$.

The exceptional length $O(\eta D_t^2)$ must be compared with the
resolution $\min(\ell,D_t^2)$ supplied by the modified denominator.
The real parts of the two Schur denominators need not vanish at the same
point. Apply \cref{imp:resolution} to $N_1$ and $s-e$, using
\cref{imp:linewidth}. It gives
\[
 \max_s W(s,t)
 \le\frac{C}{\min(\ell,D_t^2)}\int W(s,t)\,\dd s.
\]
Combining with \cref{imp:slice-length},
\[
 \int W(s,t)\one_{\mathcal E_{b'}}\,\dd s
 \le\frac{C\eta D_t^2}{\min(\ell,D_t^2)}
       \int W(s,t)\,\dd s
 \le\operatorname{poly}(L)\eta\int W(s,t)\,\dd s.
\]
The last inequality uses $D_t\le D<1$ and the inverse-polynomial lower
bound on $\ell$. Averaging over $t$ proves
\begin{equation}
 \E_{s,t}[W\one_{\mathcal E_{b'}}]
 \le\operatorname{poly}(L)\eta\,\E_{s,t}W
 \label{imp:second-case}
\end{equation}
in the small-coupling case. This estimate remains uniform as the coupling
vanishes; no positive lower bound for $D$ was used here.

\subsection{Restoring the side shift}

We now choose the parameters. All polynomial losses in
\cref{imp:comparable-case,imp:first-case,imp:second-case}, and the box
width exponent, have been fixed independently of $H$. Choose $H$ large
enough that their products with
\[
 \eta^{1/8-3\alpha_0},\qquad \eta^{1/2},\qquad\eta
\]
tend to zero, and that \cref{imp:small-b} is uniformly small in its
case. With this $H$ fixed, take $L$ sufficiently large that
$\tau\ll\Delta$, $\tau/\eta^{\alpha_0}\ll1$, and all preceding
absorptions are valid. The assumption $\tau\le e^{-c_*L}$ ensures every
one of these requirements for any fixed $c_*>0$.

For each of the eight tests, condition on the other corner variables,
apply the appropriate estimate, and compare $W$ with the same weight
$w_0$ by \cref{imp:modified-weight}. Integration and summation give
\begin{equation}
 \E[w_0\one_{\mathcal G}]\ge\frac34\E w_0,
 \qquad
 \mathcal G=\bigcap_{b'}\mathcal E_{b'}^c,
 \label{imp:weighted-Qgood}
\end{equation}
for sufficiently large $L$. This is still conditional on the independent
side uniforms. The threshold can be chosen uniformly in those frozen
values and in all exterior data. Zero conditional weights cause no
difficulty, since every estimate was an inequality between integrals.

Fix data in $\mathcal G$, before restoring the common side shift. On the
space $P$ consider the rational matrix function
\[
 F(z)=P(\mathcal T_0+zP_0)^{-1}P,
 \qquad\operatorname{Im}z<0.
\]
It is analytic there and has nonnegative imaginary part. To check
analyticity, a kernel at one point of the lower half-plane would be
annihilated by $P_0$ and by $\mathcal T_0$, hence would also be a kernel
at $z=-i$, where the inverse exists. On $\mathcal G$,
\cref{imp:Qgood} bounds the $Q,Q$ block of $F(-i)$ by $\eta^{-1}$.
Ward gives a bound one on the $P_0,P_0$ block and
\[
 \norm{P_0(\mathcal T_0-iP_0)^{-1}Q}^2
 \le\norm{Q\operatorname{Im}(\mathcal T_0-iP_0)^{-1}Q}
 \le\eta^{-1}.
\]
Transpose symmetry gives the same bound on the other mixed block. Thus
\begin{equation}
 \norm{F(-i)}\le C\eta^{-1}.
 \label{imp:center-bound}
\end{equation}

For completeness, map the unit disk to the lower half-plane by
$z=-i(1+w)/(1-w)$. Then $-iF(z(w))$ has positive semidefinite Hermitian
real part, and its value at zero obeys \cref{imp:center-bound}. The
positive-real-part weak estimate \cref{mat:positive-real-extension} gives
a boundary weak-$L^1$ bound at most
$C(\rank P)\eta^{-1}$. One can obtain this bound directly from the
positive matrix measure in its Herglotz representation: its total trace
mass is at most $(\rank P)\norm{F(-i)}$, while the constant imaginary
part has the same polynomial bound. The boundary change of variable
sends Haar measure to $\dd x/(\pi(1+x^2))$.

The Cauchy density has a positive lower bound on the fixed interval
$[-h/4,h/4]$. Therefore there is a polynomial
$R_L=C(1+L)^C$, now allowed to depend on $H$, such that for at least
half the common side shifts $x$ in their prescribed uniform law,
\begin{equation}
 \norm{P(\mathcal T_0+xP_0)^{-1}P}\le R_L.
 \label{imp:side-good}
\end{equation}
This conclusion uses only the norm at $-i$, not the norm of $T$.

Restore the artificial absorptions at the single side sites. Their
change has norm at most $\tau$ and is supported on $P_0$. For a shift
satisfying \cref{imp:side-good}, the inverse difference identity on $P$
and a Neumann series give
\[
 \norm{P\mathcal T^{-1}P}\le2R_L
\]
as soon as $\tau R_L\le1/2$. This holds for sufficiently large $L$.
Throughout the common-shift average, with or without these side
absorptions, \cref{imp:comparison} compares the weight to $w_0$ by
fixed factors. Hence at least a fixed fraction of the $\rho$-weight
survives the good shifts for every corner setting in $\mathcal G$.
Integrate this statement, use \cref{imp:weighted-Qgood}, and finally
average over the initially frozen independent side uniforms. Since
$\E\rho$ and $\E w_0$ are comparable, this proves
\cref{imp:weighted-estimate}. The pointwise comparison assertion in
\cref{imp:weighted} is already contained in \cref{imp:comparison}.
\hfill$\square$

\subsection{The Schur form used by the interior}

The following formulation makes the interface with the transmission
calculation explicit.

\begin{corollary}[Eliminating the missing directions]
\label{imp:schur-interface}
Let $N=\one-P$ on the boundary space. Almost everywhere in the boundary
law, the $N,N$ block of $\mathcal T$ is invertible. Define its Schur
matrix and transformed source by
\begin{align*}
 \widehat T
 &=P\mathcal TP-P\mathcal TN(N\mathcal TN)^{-1}N\mathcal TP,\\
 \widehat k
 &=Pk-P\mathcal TN(N\mathcal TN)^{-1}Nk.
\end{align*}
Then $\widehat T$ is complex symmetric and passive, and
\begin{align*}
 \widehat T^{-1}&=P\mathcal T^{-1}P,\\
 \widehat k&=\widehat T P\mathcal T^{-1}k,\\
 P(\mathcal T-iP)^{-1}k&=(\widehat T-i\one_P)^{-1}\widehat k.
\end{align*}
In particular $\rho=\norm{(\widehat T-i\one_P)^{-1}\widehat k}^2$,
and \cref{imp:weighted-estimate} controls $\widehat T^{-1}$ on a set
carrying a fixed fraction of this weight.
\end{corollary}

\begin{proof}
On the four-dimensional space $N$, the Hermitian real part is a fixed
real symmetric matrix plus the four common corner shifts on its
diagonal; the traceless box directions have zero diagonal entry there.
Its determinant is a nonzero polynomial, so the block is invertible
almost everywhere by passivity and the density of those shifts. The
identities are the block inverse and source elimination formulas.
Symmetry and passivity persist under Schur elimination by
\cref{mat:ward}.
\end{proof}
